\section{背景与相关工作}

\subsection{工具调用基准的演进}

工具调用评测已经从单次函数选择逐步扩展到多轮交互、环境状态变化和跨工具协作。早期的基准关注"智能体能否选对工具并填对参数"，随着智能体能力提升，评测焦点转向"能否在长程多步交互中可靠地完成任务"。本文使用的四个基准分别代表了这一演进的不同方向，它们为 RPL 研究提供了任务骨架，但都不直接评测隐私披露。

\subsubsection{$\tau$-bench：状态化多轮交互}

$\tau$-bench 针对的是已有基准的两个空白：它们只测单步、信息一次性给全的指令遵循，既没有人在回路的多轮交互，也不要求遵循领域特定策略\citep{yao2025taubench}。$\tau$-bench 以客服场景为切入点，让智能体在语言模型模拟的用户的动态对话中调用数据库 API，并遵守详细的策略文档。

$\tau$-bench 的奖励函数是二值的：动作正确性（最终数据库状态是否与标注一致）乘以输出正确性（回复是否包含用户所需信息）。在此基础上，它提出 pass$^k$ 指标衡量跨试验的一致性。两个领域各有特点：$\tau$-retail 处理订单取消、退换和地址修改，规则较简单；$\tau$-airline 涉及订票、改签和退款，策略更复杂且与会员等级和舱位挂钩，构成多跳推理难题。

$\tau$-bench 的用户档案包含 email、地址、支付方式和出生日期等字段，天然适合 RPL 研究。但其关键局限在于：奖励函数只比对最终数据库写状态，身份认证是读操作、不进入奖励，因此跳过认证也能获得满分，越权或跨用户泄露不会被惩罚\citep{yao2025taubench}。换言之，$\tau$-bench 测的是"把事做对"，而非"把事做得隐私合规"。

\subsubsection{BFCL：函数调用的综合评测}

BFCL 的动机是此前没有一个标准化的函数调用基准——判断一次调用是否"有效"本身很难，获取多样且真实的函数集合也不易\citep{patil2025bfcl}。BFCL 覆盖单轮（AST 匹配与真实执行）、多轮和智能体三类场景，跨 Python、Java、JavaScript 等多语言。

BFCL 的多轮数据定义为一个 query 序列与每轮期望的函数调用轨迹。它包含四个变体：base 是标准多轮任务；miss\_param 故意省略参数以测试追问能力；miss\_func 移除函数以测试澄清能力；long\_context 加长上下文以测试长程准确性。评测同时检查状态正确性（最终系统状态）和响应正确性（是否走了最小必要调用序列）。

BFCL 的局限与本文研究目标直接相关：它仅检查"调用是否正确"，不检查"调用是否该发生"。附录中仅提到"用户提示中的敏感信息被替换为占位符以保护隐私"\citep{patil2025bfcl}，但完全不考察模型是否应拒绝把敏感信息写进工具参数。其多轮 API codebase 是合成的模拟实现（如填油箱、查股价），不涉及真实企业级鉴权或数据最小化约束。

\subsubsection{API-Bank：工具检索与规划}

API-Bank 围绕三个递进能力构建：规划（plan）、检索（retrieve）和调用（call）\citep{li2023apibank}。它用"API 池大小"和"单轮调用次数"两个维度划分出三个 level：Level-1 评测已知 API 的参数填充，Level-2 要求先用 ToolSearcher 检索再调用，Level-3 要求自主规划并多步调用。评测集分别有 214、50 和 50 条对话。其 73 个 API 由工程师真实实现，涵盖天气、账号管理、日程、健康和财务等领域。

API-Bank 的对话大量涉及用户名、密码、token 和薪资等敏感信息，但论文未对工具调用链中的隐私风险做任何评测\citep{li2023apibank}。其伦理声明仅承诺访谈不泄露个人信息，未覆盖工具参数中的敏感数据处理。此外，Level-3 仅 50 条对话且复杂度有限，大部分对话是单步任务，多步推理能力远未被充分测试。

\subsubsection{AppWorld：可执行应用与状态测试}

AppWorld 认为已有基准只需 1--4 个 API 的线性调用，既不需要"富代码"也不需要智能体根据交互结果迭代改写代码\citep{trivedi2024appworld}。它构建了 9 个日常应用（Gmail、Venmo、Amazon、Spotify 等）的沙盒环境，包含 457 个 API 和约 37 万行底层数据库，让智能体完成跨应用的日常数字任务。

AppWorld 的评测采用基于最终数据库状态的程序化单元测试，而非逐条比对话步骤——因为复杂任务解法多样。它计算起始与结束数据库的 diff，断言期望的变更全部出现且没有附带损害（如任务没让退货却发起退货）。约 15\% 的任务是问答型，还需核对答案正确性。

AppWorld 的底层数据库是真实个人数字生活的镜像（邮件、转账记录、购物、联系人、支付卡），天然适合参数级隐私泄露评测。但其状态测试只关注任务目标和附带数据库改动，对"把不需要的敏感字段塞进可选参数"这类隐蔽泄露是盲区\citep{trivedi2024appworld}。此外，其官方 reference 调用序列只证明"任务可解"，不证明参数最小化或业务上必要。

\subsection{基于情境和用途的隐私判断}

隐私风险不能仅由字段是否敏感决定。上下文完整性理论指出，信息传递是否适当取决于信息主体、发送者、接收者、信息类型和传递原则\citep{nissenbaum2004contextual}。CONFAIDE 将这一理论用于语言模型的信息分享判断\citep{mireshghallah2024confaide}；PrivacyLens 则在行动场景中评估语言模型对隐私规范的理解，并表明模型的口头判断与实际行动可能不一致\citep{shao2024privacylens}。NIST 和 OECD 的隐私框架同样强调目的限制与数据最小化\citep{nist2017privacyengineering,nist2020privacyframework,oecd2013privacy}。

在工具调用场景中，同一信息对不同工具可能具有不同必要性。支付方式对支付工具是必要信息，但对商品查询工具并非必要；电话号码对通知工具可能必要，但不应进入内部分析备注。因此，本文把隐私标签定义为"某项具体用户信息对某个工具是否为完成当前步骤所必需"，而不是给字段分配一个全局允许或禁止状态。

\subsection{ToolPrivacyBench 与本文的关系}

\ToolPrivacyBench 首次将用途限制引入工具调用智能体的隐私审计\citep{hu2026toolprivacybench}。它把隐私泄露建模为知识库审计问题：每个 case 包含用户任务、工具集、私有原子集和字段-工具授权矩阵，智能体在 mock 后端执行后，审计器检测其工具参数中是否包含 forbidden 的私有原子。其 MT-POI 指标聚合多个子指标（字段机会归一化过披露率、严重度加权率等），在 9 个主流智能体上发现任务完成率高达 92--98\% 但 MT-POI 仍达 19--28，证明任务成功与隐私合规是两个独立维度。

\ToolPrivacyBench 的公开派生数据使用与本文相同的四个上游基准作为工作流骨架，并采用相同的筛选门槛（至少四次源调用、两个不同工具、四个敏感原子和二十个 forbidden 机会）。但其数据构造方式与本文有三个本质区别：

第一，\ToolPrivacyBench 的隐私泄露是\emph{被动观察}智能体执行行为得到的——它依赖智能体在 mock 后端执行时主动把多余信息写入参数或自由文本字段。本文则\emph{主动构造}配对调用：每一步保留上游 reference 作为不泄露参照，并构造一个仅在可选参数中增加敏感字段的泄露参照。这使得泄露是确定性的、可复现的，且与智能体行为无关。

第二，\ToolPrivacyBench 的授权标注采用 tool-purpose-first 的人工加规则方式，并对 215 个抽样 case 做了双人一致性验证（Cohen's $\kappa$ = 0.84--0.88）。本文采用确定性的 reference-argument matching 启发式规则，明确声明这不是人工业务授权金标准。

第三，\ToolPrivacyBench 关注自由文本字段（message、note、summary）中的泄露，并专门设计了 FreeTextFOR 等指标。本文聚焦 optional 结构化参数，每个泄露参数精确对应一个 entity 字段，避免把多个敏感原子塞入通用文本字段。

本文不复现 \ToolPrivacyBench 的后端模型实验，也不使用其人工一致性结果。当前工作定位为数据集构建研究，重点在于如何从四个锁定来源稳定地生成可审计的 RPL 契约，而非比较不同智能体的隐私表现。

\subsection{与智能体安全研究的区别}

RPL 属于正常任务中的非对抗性过度披露，不假设攻击者篡改提示、工具或环境。AgentDojo 等工作主要研究提示注入如何诱导智能体执行恶意动作或泄露信息\citep{debenedetti2024agentdojo}。两类问题都需要观察工具调用轨迹，但威胁模型不同：提示注入关注攻击是否改变行为，本文关注在正常任务中是否披露了超出当前用途的信息。RPL 不需要攻击者存在——只要智能体拥有完整用户档案且工具接受可选参数，过度披露就可能自然发生。
